\section{Introducción}
\label{sec:intro}

\subsection{Propósito del laboratorio}
La guía del laboratorio pide desarrollar, en dos semanas, una aplicación multiplataforma que enfrente un problema \textbf{basado en requisitos reales} y que abarque un tema de \textbf{proyección social}. La aplicación debe guardar relación con el diseño especificado en los trabajos anteriores, mostrar su desarrollo según los requerimientos planteados, basar el diseño de las interfaces en el árbol de problemas e investigar los componentes del framework relacionados con el tema del proyecto \cite{guia05}.

Para este trabajo continué con \textbf{MatrLab}, la aplicación de matrícula en laboratorios de la Escuela Profesional de Ingeniería de Sistemas que analicé en el Laboratorio~1, prototipé en Adalo en el Laboratorio~2 y programé en Flutter en el Laboratorio~4 \cite{lab1,lab2,lab4}. La diferencia es de método: en el Laboratorio~4 los requisitos salían de hipótesis mías; ahora salen de \textbf{casos documentados de matrícula en la Universidad Nacional de San Agustín (UNSA) durante 2026}, se especifican con atributos verificables, se priorizan en un \emph{Product Backlog} y se construyen en dos \emph{sprints} de Scrum adaptado a un equipo de una persona. La versión resultante es MatrLab~2.0.

\subsection{Trabajo realizado y su ubicación en el informe}
\begin{table}[H]
	\centering\small
	\begin{tabularx}{\textwidth}{|p{4.0cm}|X|p{2.4cm}|}
		\hline
		\cab\hd{Pedido de la guía} & \hd{Trabajo realizado} & \hd{Dónde}\\
		\hline
		Problema basado en requisitos reales y de proyección social & Reuní 8 evidencias públicas (prensa local, comunicados de la UNSA, INEI) sobre las matrículas de 2026 y de 2016, y las convertí en 8 hallazgos & Secc.~\ref{sec:contexto} \\\hline
		Relación con el diseño de los trabajos anteriores & Se conserva MatrLab (RF1 a RF7 del Laboratorio~4) y se amplía su árbol de problemas de 6 a 7 causas, con trazabilidad a los requisitos nuevos & Secc.~\ref{sec:requisitos} \\\hline
		Cómo plantear los requisitos & Proceso en cinco pasos, 5 requisitos funcionales nuevos con atributos, 6 no funcionales por ISO/IEC 25010, backlog con puntos y prioridad MoSCoW, criterios de aceptación Dado/Cuando/Entonces y matriz de trazabilidad & Secc.~\ref{sec:requisitos} \\\hline
		Proceso Scrum & Meta de producto, dos sprints con meta y capacidad, definición de terminado, eventos adaptados y retrospectiva & Secc.~\ref{sec:scrum} \\\hline
		Investigar componentes del framework para el tema & Estudié 12 componentes de Flutter (selección, retroalimentación, accesibilidad, estado y tiempo) y descarté tres paquetes con su razón & Secc.~\ref{sec:componentes} \\\hline
		Interfaces basadas en el árbol de problemas & Dos pestañas nuevas (Ayuda y Operación), cuatro pantallas modificadas y seis mockups con cada elemento asociado a su widget & Secc.~\ref{sec:diseno} \\\hline
		Mostrar el desarrollo según los requerimientos & Cada historia de usuario se muestra con su código, su captura y su prueba & Secc.~\ref{sec:desarrollo} \\\hline
		Mostrar el sistema & Recorrido solo con capturas: estudiante, casos especiales, administrador y tableta (30 capturas en total) & Secc.~\ref{sec:recorrido} \\\hline
		Resultados de la ejecución & Previstos para analizar el código, correr 88 pruebas y generar las capturas; quedan \textbf{pendientes} de la ejecución en mi equipo & Secc.~\ref{sec:resultados} \\\hline
	\end{tabularx}
	\caption{Resumen de lo desarrollado y su ubicación en el informe.}
	\label{tab:cumplimiento}
\end{table}

\subsection{Alcance y límites}
\begin{itemize}
	\item \textbf{Qué es:} una aplicación Flutter con 11 pantallas, lógica de matrícula con turnos por ciclo, número de atención virtual, comprobante con código y modo de contingencia, más un panel de operación para el administrador.
	\item \textbf{Qué no es:} no tiene servidor ni base de datos; los datos viven en memoria y se reinician al cerrar la aplicación. \textbf{La caída del servicio es una simulación} que activa el administrador con un interruptor; no mide ni mejora la capacidad real de un servidor. La concurrencia masiva (la causa C1 del árbol) solo se resuelve con infraestructura, y así lo reconoce el propio director de Servicios Académicos de la UNSA \cite{elpueblo25}.
	\item \textbf{Evidencia:} es documental (noticias y comunicados públicos). \textbf{No realicé encuestas ni entrevistas} a estudiantes; por eso los datos de usuarios que aparecen son los que las fuentes citan, y lo demás son decisiones de diseño que se marcan como hipótesis.
	\item \textbf{Capturas:} se generan ejecutando la aplicación con el motor de pruebas de Flutter (\texttt{flutter\_test}); no son fotos de un teléfono físico. Las 13 capturas nuevas (c18 a c30) y los resultados de las pruebas se incorporan cuando ejecuto el proyecto (sección~\ref{sec:resultados}).
	\item \textbf{Datos de ejemplo:} los cursos y grupos son los del Laboratorio~4. Los turnos de matrícula, la cola de atención de demostración y los tiempos de espera son valores de demostración míos, no datos oficiales de la UNSA.
\end{itemize}
